\section{A suggested capstone rubric}\label{appE:sec:rubric}
The table below suggests one way to assess the independent capstone of Section~\ref{ch18:sec:independent}. Each row names one quality of a finished project, gives its share of the grade, and describes what a strong project shows. The rows follow the list of contents for the write-up in that section, and the row ``Understanding'' refers to the short oral discussion that ends the investigation.
\begin{center}
\begin{tabularx}{\textwidth}{@{}p{1.15in}p{0.5in}X@{}}\toprule
Quality & Weight & What a strong project shows\\\midrule
Claim and scope & 20\% & The exact objects, hypotheses and conclusion of each result. If the question changed during the project, both the original and the revised question are stated, and the difference between them is explained. Each claim is marked with the evidence for it: proved, cited from another source, checked by computer for finitely many cases, or not established.\\
Argument & 35\% & A complete proof or a counterexample, with a map of the background that the argument needs. The delicate steps are written out in full, and each step names the results it depends on.\\
Understanding & 20\% & In the oral discussion, the student explains a delicate step without help, and, when an assumption or an example is changed, correctly predicts what still holds and what fails.\\
Reproducibility & 15\% & Computer checks where they are relevant, with the program, the range of cases it covers and the software used, so that anyone can repeat them. If no computation is relevant, the write-up explains how a reader can check the claims without a computer, for example by working small cases by hand.\\
Context & 10\% & Accurate credit to every source used; a modest statement of what is new compared with the published literature and what is new only to the student; an open account of how an AI assistant was used and of what remains unsettled.\\\bottomrule
\end{tabularx}
\end{center}
Three faults cannot be outweighed by good presentation: a statement that has been changed without saying so, a missing step on which the argument depends, and a claim of novelty without support. A project with any of these faults needs to be corrected before it can pass. The rubric measures how carefully the student carries out a small investigation of their own. It does not ask whether a famous problem was solved.
